\chapter*{Abstract}

Coherent Elastic Neutrino-Nucleus Scattering (CE$\nu$ NS) is a neutral current process proposed in 1974. In this process, a neutrino scatters off a nucleus as a whole with a low moment transfer, resulting in a very small nuclear recoil energy. Its coherent nature leads to a relatively larger cross-section than other neutrino interactions. However, the experimental measurements of CE$\nu$NS pose challenges due to the subtle recoil energy involved. In 2017, it was measured by the COHERENT Collaboration for the first time using CsI detectors and subsequently in 2020 with liquid argon (LAr) detectors. We study an array of detectors made with the same element but with a different number of neutrons, that is, different isotopes, thereby generating detectors with correlated systematic uncertainties. This correlation enhances measurement precision for testing Standard Model predictions, nuclear physics, and new physics scenarios. Our study considers two detector arrays: one using two zinc (Zn) isotopes and the other using two silicon (Si) isotopes. With these arrays, we explored the quadratic dependence of the number of neutrons in the CE$\nu$NS cross-section within different recoil energy ranges. Neutrino flux was employed from $\pi$-DAR (pion decay at rest) and reactor sources. Lastly, we use non-standard interactions and the sensitivity to the weak mixing angle as examples of the capabilities of this future detector.

